\documentclass[12pt]{article}
\usepackage[a4paper,margin=1in]{geometry}
\usepackage{amsmath,amssymb}
\usepackage{enumitem}
\usepackage{hyperref}
\usepackage{titlesec}

\titleformat{\section}{\large\bfseries}{\thesection.}{0.5em}{}
\titleformat{\subsection}{\normalsize\bfseries}{\thesubsection.}{0.5em}{}

\title{Abstract Algebra - Stockholm University 2025}
\author{}
\date{}

\begin{document}

\maketitle

\section{Groups}
\subsection{Conjugation, centralizer, center, normalizer}
For $g\in G$, the conjugation map $c_g(x)=gxg^{-1}$ is an automorphism of $G$ (an \emph{inner} automorphism).
For a subset $S\subseteq G$, the centralizer is
\[
C_G(S)=\{\,g\in G:\ gs=sg\ \text{ for all }s\in S\,\};\quad C_G(g):=C_G(\{g\}).
\]
The center is 
\[ 
    Z(G)=C_G(G)=\{g\in G:\ gx=xg\ \forall x\in G\}
\]
(always a characteristic, hence normal, subgroup).
For a subgroup $H\le G$, the normalizer is
\[
N_G(H)=\{\,g\in G:\ gHg^{-1}=H\,\},
\]
the largest subgroup of $G$ in which $H$ is normal (so $H\trianglelefteq N_G(H)$). One always has $C_G(H)\le N_G(H)$.
The number of conjugates of $H$ in $G$ equals $[G:N_G(H)]$.

\subsection{Normal subgroups and quotients}
A subgroup $N\le G$ is \textbf{normal} (write $N\trianglelefteq G$) if $gNg^{-1}=N$ for all $g\in G$.
Equivalently, every left coset equals the corresponding right coset: $gN=Ng$.
If $N\trianglelefteq G$, the set of cosets $G/N=\{gN:g\in G\}$ is a group with $(gN)(hN)=(gh)N$.


\subsection{Isomorphism theorems}
\begin{enumerate}[label=(\arabic*)]
\item \textbf{First isomorphism theorem.} For $\varphi:G\to H$ a homomorphism,
\[
G/\ker\varphi \ \cong\ \mathrm{im}\,\varphi.
\]
\item \textbf{Second isomorphism theorem.} For $A\le G$ and $N\trianglelefteq G$,
\[
A\cap N\trianglelefteq A,\qquad AN\le G,\qquad AN/N\ \cong\ A/(A\cap N).
\]
\item \textbf{Third isomorphism theorem.} If $K\le H\trianglelefteq G$, then $H/K\trianglelefteq G/K$ and
\[
(G/K)/(H/K)\ \cong\ G/H.
\]
\item \textbf{Correspondence (lattice) theorem.} For $N\trianglelefteq G$, the map
\[
\{\,\text{subgroups }H\mid N\le H\le G\,\}\ \longleftrightarrow\ \{\,\text{subgroups of }G/N\,\},\quad
H\mapsto H/N
\]
is a bijection preserving inclusion and index; it carries normal subgroups $H\trianglelefteq G$ with $N\le H$ to normal subgroups of $G/N$.
\end{enumerate}

\section{Group actions}
\section{Sylows Theorem}
\section{Rings}
\section{Integral -, Unique Factorization - \& Principle Ideal Domains}
\section{Field Extensions}

\end{document}
